\documentclass{article}

% Use NeurIPS 2015 style (can be updated to newer format later)
\usepackage{nips15submit_e}

% Packages
\usepackage{amsmath,amssymb,amsthm}
\usepackage{graphicx}
\usepackage{hyperref}
\usepackage{booktabs}
% \usepackage{algorithm}  % Not needed for this draft
% \usepackage{algorithmic}

% Theorem environments
\newtheorem{theorem}{Theorem}
\newtheorem{proposition}[theorem]{Proposition}
\newtheorem{lemma}[theorem]{Lemma}
\newtheorem{corollary}[theorem]{Corollary}
\newtheorem{definition}[theorem]{Definition}
\newtheorem{assumption}{Assumption}
\newtheorem{remark}{Remark}

% Custom commands
\newcommand{\R}{\mathbb{R}}
\newcommand{\E}{\mathbb{E}}
\newcommand{\Prob}{\mathbb{P}}
\newcommand{\cD}{\mathcal{D}}
\newcommand{\cL}{\mathcal{L}}
\newcommand{\cM}{\mathcal{M}}
\newcommand{\cS}{\mathcal{S}}

% Draft markers
\newcommand{\todo}[1]{\textcolor{red}{[TODO: #1]}}
\newcommand{\placeholder}[1]{\textcolor{blue}{[PLACEHOLDER: #1]}}
\newcommand{\citationneeded}{\textcolor{orange}{[Citation needed]}}

\usepackage{xcolor}

\title{Unifying Neural Race Reduction with Multi-View Theory:\\ A Mechanistic Account of Knowledge Distillation}

\author{
[Author Names] \\
[Affiliations] \\
\texttt{[emails]}
}

\nipsfinalcopy % Uncomment for camera-ready version

\begin{document}

\maketitle

%==============================================================================
% ABSTRACT
%==============================================================================
\begin{abstract}
Knowledge distillation enables a student network to learn from soft labels produced by a teacher, often achieving performance that exceeds what the student could learn from hard labels alone. Recent multi-view theory explains this by positing that neural networks converge to using a single ``view'' (feature subset) per class, while teacher ensembles cover multiple views. However, \emph{no existing theory predicts which view will be learned} or explains the mechanism by which soft labels transfer multi-view knowledge.

We address this gap by connecting multi-view theory to the neural race reduction framework. We show that (1) each view corresponds to a distinct computational \emph{pathway} in gated deep linear networks, (2) training dynamics create a winner-take-all \emph{race} where the pathway with highest initial advantage dominates, and (3) knowledge distillation circumvents this race because soft labels distribute gradient signal across all pathways the teacher has learned.

Our analysis yields the first quantitative prediction for which view a network will learn, determined by the product of input-output correlation strength and initial pathway magnitude. We validate our theoretical predictions on synthetic multi-view data, demonstrating that (i) single networks converge to $\sim 1/M$ view coverage, (ii) our formula predicts the winning view with $>80\%$ accuracy, and (iii) distillation from multi-view teachers enables students to match teacher coverage. \placeholder{Final experimental numbers}
\end{abstract}

%==============================================================================
% 1. INTRODUCTION
%==============================================================================
\section{Introduction}

Knowledge distillation (KD) presents a puzzle: why can a student network learn from soft labels what it could not learn from the original data directly? The student sees the same input distribution as the teacher, yet consistently achieves better generalization when trained on teacher outputs rather than ground-truth labels \cite{hinton2015distilling}.

Recent theoretical work by Allen-Zhu \& Li~\cite{allen2023towards} offers an explanation through the lens of \emph{multi-view learning}. They propose that real-world data contains multiple ``views''---distinct feature subsets that are each independently sufficient for correct classification. A key finding is that individual networks converge to using only \emph{one view per class}, while ensembles of networks collectively cover more views. Knowledge distillation then transfers this broader coverage from teacher ensemble to student.

However, multi-view theory leaves a critical question unanswered: \emph{which view will a network learn?} The theory describes \emph{what} happens (single-view convergence, ensemble coverage) but not \emph{why} it happens or how to predict it. As Allen-Zhu \& Li acknowledge, ``the mechanism by which networks select among equally predictive features remains unclear'' \citationneeded.

\paragraph{Our Contribution.} We address this gap by connecting multi-view theory to the \emph{neural race reduction} framework developed by Saxe et al.~\cite{saxe2022neural}. The neural race reduction shows that ReLU networks can be analyzed as gated deep linear networks (GDLNs), where different input-dependent gating patterns define distinct computational \emph{pathways} that race to explain the data.

Our key insight is that under multi-view data, each view naturally corresponds to a distinct pathway. This correspondence allows us to apply the rich dynamical systems analysis of neural races to explain multi-view learning. We make three main contributions:

\begin{enumerate}
    \item \textbf{View-Pathway Correspondence (Theorem~\ref{thm:correspondence}):} We formally establish that under multi-view data satisfying natural assumptions, each view $(y, m)$ maps to a distinct computational pathway $P_{y,m}$. The network's output decomposes as a sum over active view pathways.

    \item \textbf{Predictive Theory for View Selection (Theorem~\ref{thm:race}):} We show that pathway strengths follow winner-take-all dynamics. The winning view is determined by the \emph{initial advantage}:
    \begin{equation}
        A_{y,m} = \sigma_1(\Sigma_{y,m}) \cdot s_{y,m}(0)
    \end{equation}
    where $\sigma_1(\Sigma_{y,m})$ is the input-output correlation strength for view $m$ and $s_{y,m}(0)$ is the initial pathway strength. This provides the first quantitative prediction for which view will be learned.

    \item \textbf{Mechanism for KD Success (Theorem~\ref{thm:kd}):} We analyze the gradient structure under soft labels and show that KD introduces an \emph{external} gradient signal proportional to the teacher's response to each view. This external signal breaks the winner-take-all dynamics, enabling multiple pathways to coexist.
\end{enumerate}

\paragraph{Paper Organization.} Section~\ref{sec:background} reviews multi-view theory and neural race reduction. Section~\ref{sec:setup} introduces our formal setup and proves the view-pathway correspondence. Section~\ref{sec:race} analyzes race dynamics and derives our predictive formula. Section~\ref{sec:kd} explains why KD circumvents the race. Section~\ref{sec:experiments} validates our predictions experimentally. Section~\ref{sec:discussion} discusses limitations and future work.

%==============================================================================
% 2. BACKGROUND
%==============================================================================
\section{Background}
\label{sec:background}

\subsection{Multi-View Theory}
\label{sec:multiview}

Allen-Zhu \& Li~\cite{allen2023towards} formalize multi-view learning as follows. Data consists of $K$ classes, each associated with $M$ distinct ``views.'' Each view $m$ for class $y$ corresponds to a feature vector $\phi_{y,m} \in \R^d$. An input $x$ from class $y$ contains a random subset of active views:
\begin{equation}
    x = \sum_{m \in \cS} \phi_{y,m} + \varepsilon
\end{equation}
where $\cS \subseteq [M]$ is the set of active views and $\varepsilon$ is noise. Crucially, views are approximately \emph{orthogonal}: $\langle \phi_{y,m}, \phi_{y',m'} \rangle \approx 0$ for $(y,m) \neq (y',m')$, and each view is independently \emph{sufficient} for correct classification.

\paragraph{Key Results.} Allen-Zhu \& Li establish three main findings:
\begin{enumerate}
    \item \textbf{Single-view convergence:} Networks trained with hard labels converge to using approximately one view per class. Defining view coverage as the fraction of views the network can correctly classify in isolation, $C(f) \approx 1/M$ after training.

    \item \textbf{Ensemble diversity:} Different random seeds lead to different views being learned. An ensemble of $N$ independently trained networks covers approximately $M \cdot (1 - (1-1/M)^N)$ views.

    \item \textbf{Distillation transfers coverage:} A student trained via KD from a multi-view teacher achieves $C(S) \approx C(T)$, inheriting the teacher's view coverage.
\end{enumerate}

\paragraph{The Open Question.} While these results characterize the \emph{outcome} of training, they do not explain the \emph{mechanism}. What determines which view wins? Why does KD transfer coverage when the student sees the same data distribution?

\subsection{Neural Race Reduction}
\label{sec:race_background}

Saxe et al.~\cite{saxe2014exact,saxe2022neural} develop the neural race reduction framework for analyzing learning dynamics in deep networks.

\paragraph{Gated Deep Linear Networks.} Any ReLU network can be written as a \emph{gated deep linear network} (GDLN):
\begin{equation}
    f(x; W) = W_L \cdot D_{L-1}(x) \cdot W_{L-1} \cdots D_1(x) \cdot W_1 \cdot x
\end{equation}
where $D_\ell(x) = \text{diag}(G_\ell(x))$ are diagonal gating matrices with $[G_\ell(x)]_j = \mathbf{1}[\text{pre-activation}_j > 0]$. For fixed gating pattern, the function is linear in $x$.

\paragraph{Pathways and Races.} Different gating patterns define different computational \emph{pathways}. Under gradient descent, pathways ``race'' to explain the data. Saxe et al.~show that pathway strengths evolve according to winner-take-all dynamics: the pathway with highest initial correlation with the target grows fastest and eventually dominates.

\paragraph{Key Insight.} The neural race framework provides exactly the dynamical analysis missing from multi-view theory. If views correspond to pathways, then race dynamics explain single-view convergence.

%==============================================================================
% 3. SETUP AND THEOREM 1
%==============================================================================
\section{Setup and View-Pathway Correspondence}
\label{sec:setup}

\subsection{Problem Setup}

\begin{definition}[Multi-View Distribution]
A $(K, M, \mathbf{p})$-multi-view distribution $\cD$ over $\R^d \times [K]$ consists of:
\begin{itemize}
    \item $K$ classes indexed by $y \in [K]$
    \item $M$ views per class indexed by $m \in [M]$
    \item View features $\phi_{y,m} \in \R^d$ for each $(y,m)$
    \item View activation probabilities $\mathbf{p} = (p_1, \ldots, p_M)$
\end{itemize}
Sampling $(x, y) \sim \cD$: draw $y$ uniformly, draw active views $\cS \subseteq [M]$ where each $m$ is included independently with probability $p_m$, then set $x = \sum_{m \in \cS} \phi_{y,m} + \varepsilon$ with $\varepsilon \sim \mathcal{N}(0, \sigma^2 I)$.
\end{definition}

We make the following assumptions on the data distribution:

\begin{assumption}[View Sufficiency]
\label{ass:sufficiency}
Each view is independently sufficient for correct classification with margin $\gamma > 0$.
\end{assumption}

\begin{assumption}[View Orthogonality]
\label{ass:orthogonality}
Views are approximately orthogonal: $|\langle \phi_{y,m}, \phi_{y',m'} \rangle| \leq \delta \|\phi_{y,m}\| \|\phi_{y',m'}\|$ for $(y,m) \neq (y',m')$, where $\delta \geq 0$ is small.
\end{assumption}

\begin{assumption}[View Normalization]
\label{ass:normalization}
All views have comparable norm: $c_{\min} \leq \|\phi_{y,m}\| \leq c_{\max}$ for all $(y,m)$.
\end{assumption}

\begin{assumption}[View-Gating Distinguishability]
\label{ass:gating}
Different views induce different gating patterns: for $(y,m) \neq (y',m')$, there exists layer $\ell$ such that $G_\ell(\phi_{y,m}) \neq G_\ell(\phi_{y',m'})$.
\end{assumption}

\begin{assumption}[Gating Additivity]
\label{ass:additivity}
In the low-noise regime, gating under multiple views is approximately the OR of individual view gatings: $G_\ell^{(\cS)}(x) \approx \bigvee_{m \in \cS} G_\ell^{(y,m)}$.
\end{assumption}

\subsection{Pathways and View Responses}

\begin{definition}[View Pathway]
The pathway for view $(y,m)$ is the linear map $P_{y,m} : \R^d \to \R^K$ defined by:
\begin{equation}
    P_{y,m}(x) = W_L \cdot \text{diag}(G_{L-1}(\phi_{y,m})) \cdot W_{L-1} \cdots \text{diag}(G_1(\phi_{y,m})) \cdot W_1 \cdot x
\end{equation}
with associated pathway matrix $\mathbf{P}_{y,m} \in \R^{K \times d}$.
\end{definition}

\begin{definition}[View Response]
The response of network $f$ to view $(y,m)$ is:
\begin{equation}
    R_{y,m}(f) := P_{y,m}(\phi_{y,m}) = \mathbf{P}_{y,m} \cdot \phi_{y,m} \in \R^K
\end{equation}
\end{definition}

\begin{definition}[View Coverage]
The view coverage of network $f$ is:
\begin{equation}
    C(f) := \frac{1}{KM} \sum_{y=1}^{K} \sum_{m=1}^{M} \mathbf{1}\left[ \arg\max_k [R_{y,m}(f)]_k = y \right]
\end{equation}
\end{definition}

\subsection{Theorem 1: View-Pathway Correspondence}

\begin{theorem}[View-Pathway Correspondence]
\label{thm:correspondence}
Let $\cD$ be a $(K, M, \mathbf{p})$-multi-view distribution satisfying Assumptions~\ref{ass:sufficiency}--\ref{ass:additivity}. Let $f$ be a GDLN. For any sample $x \sim \cD_y^{\cS}$ (class $y$, active views $\cS$):

\textbf{(Part A --- Structural Decomposition)}
\begin{equation}
    f(x) = \sum_{m \in \cS} R_{y,m}(f) + E(x, \cS)
\end{equation}
where the error satisfies $\|E(x, \cS)\| \leq O(\delta \cdot |\cS|^2 \cdot \max_m \|R_{y,m}\| + \sigma \sqrt{d} \|P_{\max}\|_{\text{op}})$.

\textbf{(Part B --- View-Pathway Bijection)} Under Assumption~\ref{ass:gating}, the map $(y,m) \mapsto P_{y,m}$ is injective.

\textbf{(Part C --- Learning Equivalence)} View $(y,m)$ is detected by $f$ if and only if $\mathbf{P}_{y,m} \cdot \phi_{y,m}$ has its maximum component at index $y$.
\end{theorem}

\paragraph{Proof Sketch.} Part A follows from the conditional linearity of GDLNs: for fixed gating, the output is linear in input. Under Assumption~\ref{ass:additivity}, when multiple views are present, the gating is the OR of individual view gatings. The error term captures cross-view interactions (bounded by orthogonality) and noise. Part B is direct from Assumption~\ref{ass:gating}. Part C follows from the definition of view detection. \placeholder{Full proof in Appendix A}

\paragraph{Interpretation.} Theorem~\ref{thm:correspondence} establishes that learning a view is equivalent to strengthening its corresponding pathway. This reduces the multi-view learning problem to analyzing pathway dynamics.

%==============================================================================
% 4. THEOREM 2: RACE DYNAMICS
%==============================================================================
\section{Race Dynamics and View Selection}
\label{sec:race}

Building on the view-pathway correspondence, we now analyze how pathway strengths evolve during training.

\subsection{Pathway Strength Dynamics}

\begin{definition}[Pathway Strength]
The strength of pathway $(y,m)$ at time $t$ is:
\begin{equation}
    s_{y,m}(t) := \|\mathbf{P}_{y,m}(t)\|_F
\end{equation}
\end{definition}

\begin{definition}[View Correlation]
The input-output correlation for view $(y,m)$ is:
\begin{equation}
    \Sigma_{y,m} = \frac{p_m}{K} \cdot e_y \cdot \phi_{y,m}^T \in \R^{K \times d}
\end{equation}
with correlation strength $\sigma_1(\Sigma_{y,m}) = \frac{p_m}{K} \|\phi_{y,m}\|$.
\end{definition}

\begin{definition}[Initial Advantage]
The initial advantage of pathway $(y,m)$ is:
\begin{equation}
    A_{y,m} := \sigma_1(\Sigma_{y,m}) \cdot s_{y,m}(0)
\end{equation}
\end{definition}

\subsection{Theorem 2: Race Dynamics}

\begin{theorem}[Race Dynamics Under Multi-View Data]
\label{thm:race}
Let $\cD$ be a $(K, M, \mathbf{p})$-multi-view distribution satisfying Assumptions~\ref{ass:sufficiency}--\ref{ass:additivity}. Let $f$ be a GDLN with $L=2$ layers trained via gradient flow on cross-entropy loss. Initialize weights i.i.d.\ from $\mathcal{N}(0, \sigma_0^2/n)$.

\textbf{(Part A --- Pathway Dynamics)} The pathway strengths evolve according to:
\begin{equation}
    \frac{ds_{y,m}}{dt} = \sigma_1(\Sigma_{y,m}) \cdot s_{y,m} \cdot \left(1 - \frac{\sum_{m'} s_{y,m'}^2}{s_{\max}^2}\right) + O(\delta + \sigma_0^2)
\end{equation}

\textbf{(Part B --- Early Phase)} For $t \in [0, T_{\text{early}}]$:
\begin{equation}
    s_{y,m}(t) \approx s_{y,m}(0) \cdot \exp\left(\sigma_1(\Sigma_{y,m}) \cdot t\right)
\end{equation}

\textbf{(Part C --- Winner Determination)} With high probability over initialization, for each class $y$:
\begin{enumerate}
    \item There exists a unique winning view: $m^*(y) = \arg\max_{m \in [M]} A_{y,m}$
    \item The winning pathway dominates: $\lim_{t \to \infty} s_{y,m}(t) / s_{y,m^*}(t) = 0$ for $m \neq m^*(y)$
\end{enumerate}

\textbf{(Part D --- Predictive Formula)} The probability that view $m$ wins for class $y$ is:
\begin{equation}
    \Prob(m^*(y) = m) = \Prob\left( A_{y,m} > \max_{m' \neq m} A_{y,m'} \right)
\end{equation}
Under symmetric views: $\Prob(m^*(y) = m) = 1/M$.

\textbf{(Part E --- View Coverage)} The trained network satisfies $C(f) = 1/M + O(1/\sqrt{K})$.
\end{theorem}

\paragraph{Proof Sketch.} Part A follows from computing the gradient of pathway strength under the loss, decomposing by view contributions. The dynamics are of Lotka-Volterra competitive form. Part B follows from neglecting the competition term when all strengths are small. Parts C-D follow from analyzing the competitive dynamics: the pathway with highest initial advantage captures all the ``learning budget'' for its class. Part E follows from averaging over random initialization. \placeholder{Full proof in Appendix B}

\paragraph{Key Implications.}
\begin{itemize}
    \item \textbf{Initialization determines winner:} Under symmetric views (equal $\sigma_1$), the view that wins is determined purely by random initialization---whichever pathway happens to be strongest at $t=0$.

    \item \textbf{Explains single-view convergence:} The winner-take-all dynamics explain why $C(f) \approx 1/M$: each class converges to one dominant view.

    \item \textbf{Explains ensemble diversity:} Different seeds $\Rightarrow$ different initial $s_{y,m}(0)$ $\Rightarrow$ different winners.
\end{itemize}

%==============================================================================
% 5. THEOREM 3: WHY KD WORKS
%==============================================================================
\section{Why Knowledge Distillation Circumvents the Race}
\label{sec:kd}

We now analyze how soft labels from a teacher change the learning dynamics.

\subsection{The Puzzle}

From Theorem~\ref{thm:race}, hard label training leads to winner-take-all dynamics and single-view convergence. Yet empirically, students trained via KD from multi-view teachers achieve $C(S) \approx C(T)$.

The student sees the \emph{same data distribution} as the teacher. What property of soft labels enables multi-view learning?

\subsection{Gradient Analysis}

\paragraph{Hard Labels.} Under cross-entropy loss with hard labels, the gradient for pathway $(y,m)$ is approximately:
\begin{equation}
    \nabla_{P_{y,m}} \cL_{\text{hard}} \propto -(1 - p_y(x)) \cdot \phi_{y,m}
\end{equation}
This gradient is dominated by \emph{self-reinforcement}: strong pathways receive more gradient (through their contribution to $p_y$), creating winner-take-all dynamics.

\paragraph{Soft Labels.} Under KD loss $\cL_{\text{KD}} = \tau^2 \cdot KL(p_T^\tau \| p_S^\tau)$:
\begin{equation}
    \nabla_{P_{y,m}^S} \cL_{\text{KD}} = \underbrace{\alpha_m(T, \tau) \cdot \nabla_m^{\text{ext}}}_{\text{external (teacher)}} + \underbrace{\beta(s_{y,m}^S) \cdot \nabla_m^{\text{self}}}_{\text{internal (self)}}
\end{equation}

The \textbf{external signal} $\alpha_m(T, \tau) \propto \|R_{y,m}^T\|^2 / \tau$ depends on the teacher's response to view $m$---not on the student's current state. This provides gradient to all pathways the teacher has learned, regardless of their current strength in the student.

\subsection{Theorem 3: KD Breaks the Race}

\begin{theorem}[KD Circumvents Race Dynamics]
\label{thm:kd}
Let $T$ be a teacher with view coverage $\cM_T(y) := \{m : \|R_{y,m}^T\| > \epsilon\}$. Let $S$ be a student trained via KD at temperature $\tau$.

\textbf{(Part A --- Gradient Decomposition)} The gradient for student pathway $(y,m)$ is:
\begin{equation}
    \nabla_{P_{y,m}^S} \cL_{\text{KD}} = \alpha_m(T, \tau) \cdot \nabla^{\text{ext}}_{y,m} + \beta_m(S) \cdot \nabla^{\text{self}}_{y,m}
\end{equation}
where $\alpha_m(T, \tau) = \Theta(\|R_{y,m}^T\|^2 / (\tau K))$ and $\beta_m(S) = \Theta(\gamma \sigma_1 s_{y,m}^S)$ with $\gamma < 1$.

\textbf{(Part B --- Modified Dynamics)} Student pathway strengths evolve as:
\begin{equation}
    \frac{ds_{y,m}^S}{dt} = \left(\alpha_m(T, \tau) + \gamma \sigma_1(\Sigma_{y,m}) s_{y,m}^S \right) \cdot \left(1 - \frac{\|\mathbf{s}_y^S\|^2}{s_{\max}^2}\right)
\end{equation}

\textbf{(Part C --- Breaking Winner-Take-All)} If $|\cM_T(y)| > 1$, the system has a unique stable equilibrium with multiple non-zero pathway strengths:
\begin{equation}
    s_{y,m}^S(\infty) \propto \sqrt{\alpha_m(T, \tau)} \propto \|R_{y,m}^T\| / \sqrt{\tau}
\end{equation}

\textbf{(Part D --- Coverage Inheritance)} The student's coverage satisfies:
\begin{equation}
    C(S) \geq \frac{1}{KM} \sum_{y=1}^K |\cM_T(y)| - O(1/\sqrt{\tau})
\end{equation}
\end{theorem}

\paragraph{Proof Sketch.} Part A follows from computing the KD gradient and decomposing by contribution type. The external term arises from the mismatch between student and teacher outputs when view $m$ is active. Part B follows from translating gradients into pathway strength dynamics. Part C follows from equilibrium analysis: the external signal $\alpha_m$ provides a ``floor'' that prevents weak pathways from dying. Part D follows from counting views with $\alpha_m > 0$. \placeholder{Full proof in Appendix C}

\subsection{Special Cases}

\paragraph{Single-View Teacher.} If teacher learned only view $m^*$: $\cM_T(y) = \{m^*\}$, $\alpha_m = 0$ for $m \neq m^*$. The dynamics revert to Theorem~\ref{thm:race}; KD provides no benefit over hard labels.

\paragraph{Ensemble Teacher.} An ensemble of $N$ teachers covers $|\cM_T^{\text{ens}}(y)| \approx M \cdot (1 - (1-1/M)^N)$ views. With $N \geq M$, all views are covered, and the student learns all views.

\paragraph{Self-Distillation.} Teacher learned view $m_1$; student's initialization favors $m_2$. Under KD: external signal for $m_1$, internal advantage for $m_2$. Result: student learns \emph{both} views, achieving $C(S) > C(T)$.

\subsection{The Mechanism in Summary}

\begin{table}[h]
\centering
\caption{Comparison of gradient structure and dynamics under hard vs.\ soft labels.}
\label{tab:comparison}
\begin{tabular}{@{}lll@{}}
\toprule
& \textbf{Hard Labels} & \textbf{Soft Labels (KD)} \\
\midrule
Target & $e_y$ (one-hot) & $p_T^\tau(x)$ (soft distribution) \\
Gradient for pathway $m$ & $\propto s_{y,m}$ & $\propto \alpha_m(T) + \gamma s_{y,m}$ \\
Dynamics & $\dot{s} = \sigma_1 s (1 - \text{comp})$ & $\dot{s} = (\alpha + \gamma \sigma_1 s)(1 - \text{comp})$ \\
Equilibrium & One pathway dominates & Multiple pathways coexist \\
Coverage & $\approx 1/M$ & $\approx C(T)$ \\
\bottomrule
\end{tabular}
\end{table}

The key insight is that soft labels encode \emph{which view is present} in each input, through the teacher's view-specific responses. This information creates an external gradient signal that all pathways receive, breaking the self-reinforcing winner-take-all dynamics.

%==============================================================================
% 6. EXPERIMENTS
%==============================================================================
\section{Experiments}
\label{sec:experiments}

We validate our theoretical predictions on synthetic multi-view data. \placeholder{Real data experiments if time permits}

\subsection{Synthetic Data Setup}

\begin{figure}[h]
\centering
\placeholder{Figure 1: Synthetic multi-view data visualization}
\caption{Synthetic multi-view data setup. $K=10$ classes, $M=3$ views per class. Each view occupies disjoint dimensions (``slots''), ensuring orthogonality.}
\label{fig:setup}
\end{figure}

We construct synthetic data following the multi-view model:
\begin{itemize}
    \item $K = 10$ classes, $M = 3$ views per class
    \item Input dimension $d = 300$ (100 dimensions per view)
    \item Views are exactly orthogonal (occupying disjoint coordinate subsets)
    \item View activation probability $p_m = 0.5$ for all $m$
    \item Noise level $\sigma = 0.1$
\end{itemize}

We train 2-layer ReLU networks with hidden dimension $n = 256$ using SGD.

\subsection{Experiment 1: Single-View Convergence}

\paragraph{Setup.} Train 30 networks with different random seeds using hard labels.

\paragraph{Prediction.} $C(f) \approx 1/M = 0.33$ (one view per class).

\begin{table}[h]
\centering
\caption{View coverage after hard-label training across 30 seeds.}
\label{tab:coverage}
\begin{tabular}{@{}lcc@{}}
\toprule
Metric & Prediction & Observed \\
\midrule
Mean coverage $C(f)$ & 0.33 & \placeholder{0.XX $\pm$ 0.XX} \\
Per-class views learned & 1.0 & \placeholder{X.XX $\pm$ X.XX} \\
\bottomrule
\end{tabular}
\end{table}

\subsection{Experiment 2: Winner Prediction}

\paragraph{Setup.} For each trained network, compare the predicted winning view (highest $A_{y,m}$ at initialization) against the actual learned view.

\paragraph{Prediction.} $>80\%$ agreement between predicted and actual winner.

\begin{table}[h]
\centering
\caption{Winner prediction accuracy.}
\label{tab:prediction}
\begin{tabular}{@{}lc@{}}
\toprule
Metric & Value \\
\midrule
Prediction accuracy & \placeholder{XX.X\%} \\
Correlation($A_{y,m}$, final $s_{y,m}$) & \placeholder{0.XX} \\
\bottomrule
\end{tabular}
\end{table}

\subsection{Experiment 3: Race Dynamics Visualization}

\begin{figure}[h]
\centering
\placeholder{Figure 2: Pathway strength dynamics during training}
\caption{Pathway strengths $s_{y,m}(t)$ during training. \textbf{Left:} Hard labels---one pathway dominates, others decay. \textbf{Right:} KD from ensemble---multiple pathways coexist.}
\label{fig:dynamics}
\end{figure}

\paragraph{Setup.} Track pathway strengths throughout training under (a) hard labels and (b) KD from 5-teacher ensemble.

\paragraph{Prediction.} Under hard labels: winner-take-all dynamics. Under KD: multiple pathways stabilize.

\subsection{Experiment 4: KD Coverage Transfer}

\paragraph{Setup.}
\begin{enumerate}
    \item Train ensemble of 5 teachers with hard labels
    \item Measure ensemble coverage $C(T_{\text{ens}})$
    \item Train student with hard labels: measure $C(S_{\text{hard}})$
    \item Train student with KD from ensemble: measure $C(S_{\text{KD}})$
\end{enumerate}

\paragraph{Prediction.} $C(S_{\text{KD}}) \approx C(T_{\text{ens}}) \gg C(S_{\text{hard}})$.

\begin{table}[h]
\centering
\caption{Coverage transfer via knowledge distillation.}
\label{tab:kd}
\begin{tabular}{@{}lcc@{}}
\toprule
Model & Prediction & Observed \\
\midrule
Individual teacher $C(T_i)$ & 0.33 & \placeholder{0.XX} \\
Ensemble $C(T_{\text{ens}})$ & $\sim 0.87$ & \placeholder{0.XX} \\
Student (hard labels) $C(S_{\text{hard}})$ & 0.33 & \placeholder{0.XX} \\
Student (KD) $C(S_{\text{KD}})$ & $\sim 0.87$ & \placeholder{0.XX} \\
\bottomrule
\end{tabular}
\end{table}

\subsection{Experiment 5: Gradient Distribution}

\paragraph{Setup.} Measure gradient magnitude for each pathway under hard labels vs.\ KD.

\paragraph{Prediction.} Under KD, gradient is more uniformly distributed across pathways.

\begin{figure}[h]
\centering
\placeholder{Figure 3: Gradient distribution across pathways}
\caption{Gradient magnitude per pathway during training. \textbf{Left:} Hard labels---gradient concentrated on dominant pathway. \textbf{Right:} KD---gradient distributed across teacher's views.}
\label{fig:gradient}
\end{figure}

\todo{Add real data experiments on CIFAR-10/100 if time permits}

%==============================================================================
% 7. DISCUSSION
%==============================================================================
\section{Discussion}
\label{sec:discussion}

\subsection{Limitations}

Our analysis relies on several simplifying assumptions:

\begin{itemize}
    \item \textbf{GDLN approximation:} We analyze gated deep linear networks as a proxy for ReLU networks. While the GDLN equivalence is exact~\cite{jarvis2025relu}, our dynamics analysis assumes fixed gating patterns.

    \item \textbf{Orthogonal views:} We assume views are approximately orthogonal. Real-world features may have correlations that violate this assumption.

    \item \textbf{Two-layer networks:} Our detailed dynamics analysis is for $L=2$ layers. Extension to deeper networks introduces additional complexity.

    \item \textbf{Unknown view structure:} In real data, view structure is not known. Our synthetic experiments validate the theory; real-data validation requires proxy measures.
\end{itemize}

\subsection{Future Work}

\begin{itemize}
    \item \textbf{Optimal teacher design:} Given student capacity, what teacher coverage maximizes student performance?

    \item \textbf{Progressive distillation:} Can iterative self-distillation progressively increase view coverage?

    \item \textbf{Feature distillation:} Does intermediate-layer matching provide additional view signal?

    \item \textbf{Beyond orthogonal views:} How do our results extend when views are correlated?
\end{itemize}

\subsection{Conclusion}

We have unified multi-view theory with neural race reduction to provide a mechanistic account of knowledge distillation. Our framework:

\begin{enumerate}
    \item Explains \emph{why} networks learn one view (race dynamics)
    \item Predicts \emph{which} view will be learned (initial advantage formula)
    \item Explains \emph{why} KD enables multi-view learning (external gradient signal)
\end{enumerate}

This is the first predictive theory for view selection in neural networks, filling a gap explicitly acknowledged in prior work. Our analysis suggests that the ``dark knowledge'' in soft labels is precisely the information about which view is present---information that hard labels cannot convey.

%==============================================================================
% REFERENCES
%==============================================================================
\bibliographystyle{plain}

\begin{thebibliography}{10}

\bibitem{allen2023towards}
Zeyuan Allen-Zhu and Yuanzhi Li.
\newblock Towards understanding ensemble, knowledge distillation and self-distillation in deep learning.
\newblock In \emph{International Conference on Learning Representations (ICLR)}, 2023.

\bibitem{hinton2015distilling}
Geoffrey Hinton, Oriol Vinyals, and Jeff Dean.
\newblock Distilling the knowledge in a neural network.
\newblock \emph{arXiv preprint arXiv:1503.02531}, 2015.

\bibitem{jarvis2025relu}
Devon Jarvis, Richard Klein, and Benjamin Rosman.
\newblock On the relationship between ReLU networks and gated linear networks.
\newblock In \emph{International Conference on Machine Learning (ICML)}, 2025.
\newblock \citationneeded

\bibitem{saxe2014exact}
Andrew~M. Saxe, James~L. McClelland, and Surya Ganguli.
\newblock Exact solutions to the nonlinear dynamics of learning in deep linear networks.
\newblock In \emph{International Conference on Learning Representations (ICLR)}, 2014.

\bibitem{saxe2022neural}
Andrew~M. Saxe, Shagun Sodhani, and Sam Gershman.
\newblock The neural race reduction: Dynamics of abstraction in gated networks.
\newblock In \emph{International Conference on Machine Learning (ICML)}, 2022.

\end{thebibliography}

%==============================================================================
% APPENDICES (Placeholders)
%==============================================================================
\appendix

\section{Proof of Theorem 1}
\label{app:thm1}
\placeholder{Full proof to be added when finalized}

\section{Proof of Theorem 2}
\label{app:thm2}
\placeholder{Full proof to be added when finalized}

\section{Proof of Theorem 3}
\label{app:thm3}
\placeholder{Full proof to be added when finalized}

\section{Additional Experimental Details}
\label{app:experiments}
\placeholder{Hyperparameters, training details, additional results}

\end{document}
